\section{The outward envelope and the low-frequency barrier}
\label{sec:excess}

We continue the induction for \Cref{sx:signed} at the integer \(Q\).
Thus \Cref{ct:separated} is available when \(Q>1\), using only the
assertion at smaller integers. We shall construct a positive envelope
\(I(s)\) of future signed masses \(A(s)=s^2M_s\) and follow its
logarithmic level \(a(s)=\log(s^2/I(s))\). The output in
\Cref{ex:bands} is a traversal from \(a=B\) to \(a=2B\) in one of
\(Q\) fixed bands, with no merger of its height centers and with
\(k\ge B^{-\eps_0}\) for the coefficient \(k\) of the height moment,
where \(\eps_0>0\) is fixed and small. This holds for every contact
selected below whose outward vertical contains a threshold equality; \Cref{sec:packing}
will turn the traversal into many scales with normal flux.

All heights in this section are in the scaled coordinates of
\Cref{sx:signed}; the quantities defining
\(\Lambda\) retain their physical units. Gaussian averages are centered
at the contact under consideration. The dot denotes differentiation
with respect to \(t=\log s\).

Fix, in this order, a sufficiently large Gaussian exponent \(K>2m+100\),
a sufficiently large inverse-polynomial exponent \(P_0\), and the
constants in the separated estimate. They may depend on \(m,n,Q\) and
the fixed coordinate bounds. Increasing these fixed constants only
increases the lower threshold for \(T\). Put
\begin{equation}
 F=F_Q(T),\qquad \eta=e^{-2F},\qquad
 a_0=F^2,\qquad d_*=e^{-a_0},\qquad \rho=\eta^4T^{-2}.
 \label{ex:parameters}
\end{equation}
The elementary rate comparisons proved in \Cref{ex:bands} include
\(F=T^{o(1)}\), \(a_0=o(T)\), and \(|\log\rho|=o(T)\).

\subsection{Selection and the bottom boundary}

Extend the localized measures to the fixed torus of
\Cref{ct:localization}, and write \(\nu=\tau+d\,\dd y\), with
\(\tau\) restricted to a fixed interior enlargement of the working
set. For a weight \(w\) in \Cref{sx:signed}, use the source
\begin{equation}
 f=w(L-C_Q\Lambda)-\eta\nu.
 \label{ex:source}
\end{equation}
To prove the required upper bound it is enough to exclude
\(\int f\ge\eta d\): the total mass of \(\nu\) is at most \(Cd\),
and \((C+1)\eta\le e^{-F}\) for large \(T\).

Start at \(\rho\), continue downwards while \(\tau_s(x)\le d_*\),
and stop at a threshold equality or at a positive bottom \(\sigma\).
Use the continuous entrance boundary, including its vertical sides,
from \Cref{ct:contact}. At the top the continuation inequality is
strict, since \(\tau_\rho\le Cd\rho^{-m}=o(d_*)\).
Let \(K_\sigma\) denote this compact entrance boundary. Write
\(\widehat h_\sigma\) for its contact measure,
\(h_\sigma=(\pi_X)_\#\widehat h_\sigma\) for the spatial projection,
and \(\mathcal S_\sigma=\spt\widehat h_\sigma\). Both measures have
mass \(\int f\). In this section a contact means a point of
\(\mathcal S_\sigma\), where \Cref{ct:contact} supplies the pointwise
and cumulative contact inequalities. At every contact
\((x,\underline s)\),
\begin{equation}
 (h_\sigma)_{\underline s}(x)\le f_{\underline s}(x)
 \le Cd_*,\qquad
 \tau_v(x)\le d_*\quad(\underline s\le v\le\rho).
 \label{ex:threshold-control}
\end{equation}
Here and below source localization first places the contact in the
fixed interior region and gives \(w(x)>0\). The constants in
\eqref{ex:threshold-control} are independent of \(\sigma\).

Define the threshold event
\[
 \mathcal G_\sigma=
 \{(x,u)\in K_\sigma:\ \tau_v(x)=d_*
                 \text{ for some }v\in[u,\rho]\}.
\]
The corresponding set of \(((x,u),v)\) is closed in
\(K_\sigma\times[\sigma,\rho]\), by continuity of the Gaussian
averages. Its projection \(\mathcal G_\sigma\) is therefore compact
and Borel. The strict top inequality excludes \(v=\rho\). If a
boundary point above the bottom had no equality, the continuation
inequalities would be uniformly strict on its compact outward
interval; continuity would then make it an interior point. Hence
\(K_\sigma\setminus\mathcal G_\sigma\) consists precisely of bottom
points with no threshold equality.

\begin{lemma}[Vanishing contribution from the bottom]
\label{ex:bottom}
The mass \(\widehat h_\sigma(K_\sigma\setminus\mathcal G_\sigma)\)
is at most
\[
 o_{\sigma\downarrow0}(1)+Cd_*d.
\]
Every contact in \(\mathcal G_\sigma\) has a scale
\(s_0\in[\underline s,\rho)\) with
\(\tau_{s_0}(x)=d_*\). The inequalities
\eqref{ex:threshold-control} hold also when
\(\underline s<s_0\).
\end{lemma}

\begin{proof}
Only the first assertion needs proof. Put
\[
 \beta_\sigma=(\pi_X)_\#
       \bigl(\widehat h_\sigma\!\restriction_{K_\sigma\setminus\mathcal G_\sigma}\bigr)
       \le h_\sigma,
\]
and let \(\mathcal C_\sigma\) be the centers of contacts in
\(K_\sigma\setminus\mathcal G_\sigma\). Assign these contacts to grid
cubes of side \(\sigma\). If such a cube \(J\) contains a contact
\(x_J\), the Gaussian lower bound on \(J\), together with the contact
inequality, gives
\[
 h_\sigma(J)\le C\sigma^m(f_\sigma(x_J))_+
 \le Cd_*\sigma^m.
\]
Since \(\beta_\sigma\le h_\sigma\), this also bounds
\(\beta_\sigma(J)\). On an occupied half-open grid cube define
\[
 g_\sigma|_J=C\sup_{x\in\mathcal C_\sigma\cap J}
                   (f_\sigma(x))_+,
 \qquad
 \overline\beta_\sigma
   =\sum_J\frac{\beta_\sigma(J)}{|J|}\mathbf1_J\,\dd y,
\]
and put \(g_\sigma=0\) on unoccupied cubes. Then
\(\overline\beta_\sigma\le g_\sigma\,\dd y\) and
\(0\le g_\sigma\le Cd_*\). For every continuous test \(\varphi\),
\[
 \left|\int\varphi\,\dd\beta_\sigma
            -\int\varphi\,\dd\overline\beta_\sigma\right|
 \le\omega_\varphi(C\sigma)\beta_\sigma(\text{chart})\longrightarrow0,
\]
where \(\omega_\varphi\) is its modulus of continuity.

For a finite measure \(\lambda\), at almost every Lebesgue point of
its absolutely continuous density one has
\[
 \sup_{|x-y|\le C\sigma}
 \big|(G_\sigma*\lambda)(x)-\lambda_{\mathrm{ac}}(y)\big|
 \longrightarrow0.
\]
For the singular part the corresponding limit is zero almost
everywhere. To check the assertion, decompose into dyadic annuli about
\(y\); the differentiation bound on the mass in each sufficiently
small ball is uniform in the indicated displacement, and the remaining
annuli have an exponentially small Gaussian weight. Apply this first
to smooth approximations of the absolutely continuous part and then
to its \(L^1\) remainder. Apply this assertion simultaneously to
\(\tau\) and \(f\). If \(\tau_{\mathrm{ac}}(y)>d_*\), the cube
containing \(y\) is eventually unoccupied, since all its contacts
satisfy \(\tau_\sigma\le d_*\). At every other simultaneous
differentiation point the uniform displacement bound controls the
supremum defining \(g_\sigma\). Therefore
\[
 \limsup_{\sigma\downarrow0}g_\sigma(y)
 \le C\mathbf1_{\{\tau_{\mathrm{ac}}\le d_*\}}
                         (f_{\mathrm{ac}})_+(y)
 \quad\text{almost everywhere}.
\]
Reverse Fatou, using the common majorant \(Cd_*\) on the fixed
chart, and the weak-test estimate above show that every bottom limit
has density bounded by
\begin{equation}
 C\,\mathbf 1_{\{\tau_{\mathrm{ac}}\le Cd_*\}}
 (f_{\mathrm{ac}})_+.
 \label{ex:bottom-density}
\end{equation}
The value of the fixed constant in the threshold is immaterial.

At a point in this set, the area formula and \Cref{lt:jacobian}
give exactly \(Q\) projection entries, counted with their integer
multiplicities. Indeed the difference between their integer count and
\(Q\) is bounded by \(C\tau_{\mathrm{ac}}<1\).
Every entry has \(e\le Cd_*\); an entry with larger tilt would itself
contribute more to the projected density of \(D\). The part on which
the projection has deficient rank projects to a Lebesgue null set,
by the area formula for rectifiable sets.
The expansion of the projection Jacobian at each remaining entry is
\[
 2(1-J_\pi)-e
 \le Ce^2+C\bigl(|z_{\mathrm{phys}}|^2+|H_S|^2\bigr).
\]
Sum this inequality with weight \(\theta/J_\pi\). Since the count is
\(Q\), its left side is exactly the density of \(L\), with respect to
base volume. Taking \(C_Q\) large enough absorbs the displayed
geometry term and gives
\((f_{\mathrm{ac}})_+\le Cd_*D_{\mathrm{ac}}\) on the set in
\eqref{ex:bottom-density}. Its integral is at most \(Cd_*d\).
Compactness and \eqref{ex:bottom-density} prove the stated uniform
limsup bound, hence the assertion.
\end{proof}

\subsection{The future envelope}

Fix a contact in \(\mathcal G_\sigma\), and suppress \(x\) and
\(\sigma\) from the notation. Define
\begin{equation}
 \begin{split}
 f_*&=L-C_Q\Lambda-\frac\eta4(\tau+d\,\dd y),\qquad
 A(s)=s^2M_s,\\
 Z(s)&=\int_{\underline s}^{s}v^2(f_*)_v\,\dd\log v,
 \qquad
 J(s)=C_Qs^2\Lambda_s+\frac\eta4s^2(\tau_s+d).
 \end{split}
 \label{ex:path-variables}
\end{equation}
The localization inequality and the contact inequalities imply
\begin{equation}
 Z(s)\ge0,\qquad (f_*)_{\underline s}\ge0,\qquad
 A(s)\ge Z(s),\qquad \dot A\ge\dot Z.
 \label{ex:bank}
\end{equation}
Here \(A(\underline s)\ge0\). To verify the last two assertions, use
\Cref{lt:gaussian} and integrate the nonnegative radial term and the
margin. The geometric error is absorbed in, for example, half of the
\(\eta\nu/4\) margin. Uniformly down to the bottom, split its
Gaussian integral at \(|y-x|/s=T\). On the body its coefficient is
\(Ce^{-\kappa T}(1+T^2)\ll\eta\); on the tail use
\eqref{ex:threshold-control} at larger Gaussian scales and the total
mass above \(\rho\). The tail is \(o(\eta d)\).
Thus this absorption uses no lower bound on \(\underline s\).
Only the first contact scale has a pointwise sign; \eqref{ex:bank}
is the sign information used at subsequent scales.

To compare configurations at different scales, we need one quantity
that controls the unsigned excess averages at larger scales up to
\(\rho\), allowing a fixed power of the radius ratio.
The next lemma obtains it from future signed mass. At a maximizing
scale \(v\) for a weighted future \(\tau\)-average, a value of \(A(v)\)
that is a sufficiently small multiple of \(v^2\tau_v\) would make
\(f_*\) negative on a short outward interval, spending more of \(Z\)
than \(0\le Z\le A\) allows. The resulting envelope also controls
the Gaussian tails needed for the height comparison.

\begin{lemma}[Future envelope]
\label{ex:envelope}
Set
\begin{equation}
 I(s)=\max\left\{
   \sup_{s\le v\le\rho}(s/v)^K A(v),\,
   (s/\rho)^K\rho^2d\rho^{-2m-4}\right\},
 \qquad a(s)=\log\frac{s^2}{I(s)}.
 \label{ex:future}
\end{equation}
Starting at a threshold scale \(s_0\), stop at the first scale
\(s_*\) with \(a(s_*)=T/2\). This scale exists, and on the path
\([s_0,s_*]\) one has
\begin{equation}
 a(s_0)=a_0+O(1),\qquad a(s)\ge a_0-O(1),\qquad
 v^2\tau_v\le C(v/s)^K I(s)\quad(s\le v\le\rho).
 \label{ex:dilations}
\end{equation}
The functions \(I,a\) are locally Lipschitz in \(\log s\). Almost
everywhere they satisfy
\begin{equation}
 \begin{cases}
 I=A,\quad \dot I=\dot A,&\text{on the reset set }\{I=A\},\\
 \dot I=KI,&\text{on its complement},
 \end{cases}
 \qquad
 |\dot a|+\dfrac{|\dot I|+|\dot Z|}{I}\le C.
 \label{ex:derivatives}
\end{equation}
Moreover \(\dot Z/I\) is locally Lipschitz, with derivative bounded
by \(C\). For every fixed \(P\), weighted scalar control gives
\begin{equation}
 |M|_s\le C(D_s+\Lambda_s)+a^{-P}I(s)/s^2.
 \label{ex:weighted-scalar}
\end{equation}
All constants are uniform in the contact and the bottom.
\end{lemma}

\begin{proof}
Write \(I_{\mathrm{end}}\) for the artificial term in
\eqref{ex:future}. The definition makes \(I\) positive and continuous.
Let
\[
 J_\tau(s)=\sup_{s\le v\le\rho}(s/v)^K v^2\tau_v.
\]
We apply the following argument first at \(s=s_0\), and then after
initialization up to the first crossing of \(a=T/2\). At \(s_0\)
the threshold and \(K>2\) give the exact value
\(J_\tau(s_0)=s_0^2d_*\), realized at \(v=s_0\), whereas
\(I_{\mathrm{end}}(s_0)=o(s_0^2d_*)\).
At any later scale before that crossing,
\(I(s)/s^2\ge e^{-T/2}\) and
\(I_{\mathrm{end}}(s)/s^2\le e^{-T+o(T)}\), so \(I\) is
realized by an actual future value of \(A\). Since \(\tau\ge|M|\),
\(J_\tau\ge I>I_{\mathrm{end}}\). At a maximizing scale \(v\),
writing \(\delta=\tau_v\), these observations give
\begin{equation}
 e^{-T/2}\le\delta\le d_*,\qquad
 \tau_t\le(t/v)^{K-2}\delta\quad(v\le t\le\rho).
 \label{ex:maximizer-bounds}
\end{equation}
The first lower bound holds at initialization as well. For any fixed
\(N\), a maximizer with \(v\ge\rho/T^N\) would, by the root
total mass bound, satisfy
\[
 \frac{J_\tau(s)}{I_{\mathrm{end}}(s)}
 \le C\rho^{m+4}T^{N(K+m-2)}=o(1),
\]
which is impossible. Thus all fixed polynomial dilations needed
below, including the body \(C\sqrt{|\log\delta|}\,v\), lie
within \(\rho\).

Here is the raw density estimate on that body. A cell centered at
\(y\), with \(|y-x|\le Rv\) and radius
\(r=c v/(1+R)\), lies with its fixed enlargement inside a ball
of radius \(C(1+R)v\) about \(x\). The Gaussian lower bound at
that radius and \eqref{ex:maximizer-bounds} imply
\begin{equation}
 r^{-m}\tau(B_{4r}(y))
 \le C(1+R)^{2m+K-2}\delta.
 \label{ex:cell-density}
\end{equation}
For \(R\le C\sqrt{|\log\delta|}\) this is small, uniformly in
\(\delta\le d_*\). Moreover \(\tau\) includes
\(|M|=|\pi_\#\mu-Q\vol_S|\), so the normalized smooth mass
average on each cell is close to that same \(Q\). The integer in
\Cref{lt:scalar} is therefore \(Q\) on every cell. The Gaussian
has a bounded ratio on a fixed enlargement of such a cell, and
\Cref{lt:weightedballs} sums the estimates with bounded overlap.

Annular summation of \eqref{ex:maximizer-bounds} bounds the Gaussian
tail beyond \(Rv\), including any fixed derivative polynomial, by
\(C\delta\operatorname{poly}(R)e^{-cR^2}\) up to \(\rho\).
Above \(\rho\) the root mass bound gives a remainder at most
\(Cd v^{-m}\operatorname{poly}(\rho/v)e^{-c(\rho/v)^2}\).
Since \(v<\rho/T^N\), \(\delta\ge e^{-T/2}\), and
\(|\log\rho|=o(T)\), this is \(o(\delta)\), uniformly even
as \(v\) tends to zero. Increasing the fixed body constant makes
the full tail \(o(\delta)\). Consequently
\[
 |M|_v\le C(D_v+\Lambda_v)+o(\tau_v),\qquad
 \tau_v\le C(D_v+\Lambda_v).
\]

Suppose \(A(v)\le\varepsilon v^2\tau_v\). For \(C_Q\) sufficiently
large and then \(\varepsilon\) sufficiently small,
\[
 (f_*)_v\le-c\tau_v.
\]
Gaussian differentiation and the maximizing dilation bound show
that the derivative, in logarithmic scale, of
\(2M_r-D_r-C_Q\Lambda_r\) is bounded by \(C_K\tau_v\) on a fixed
short outward interval. Thus, for a fixed \(\vartheta>0\),
\((f_*)_r\le-c\tau_v/2\) on \([v,e^\vartheta v]\).
The negative margin may simply be omitted in obtaining this upper
bound. By \eqref{ex:bank},
\[
 Z(e^\vartheta v)
 \le Z(v)-c'\vartheta v^2\tau_v
 \le \varepsilon v^2\tau_v-c'\vartheta v^2\tau_v<0
\]
if \(\varepsilon<c'\vartheta\). This contradicts the cumulative
contact sign. Hence \(A(v)\ge c v^2\tau_v\) at an uncontrolled
maximizer. At \(s_0\) this first proves
\(I(s_0)\ge c s_0^2d_*\); at each subsequent path scale it
proves the last inequality in \eqref{ex:dilations}. This order uses
only maximality and the root mass bound when establishing the raw
cell estimates. The same cell argument, now using the proved
dilation bound, gives \eqref{ex:weighted-scalar} at every path
scale; any prescribed inverse power is obtained by enlarging the
Gaussian body constant.

At \(s_0\), the threshold equality and the just proved bound give
\(I(s_0)\ge c s_0^2d_*\). The opposite bound follows from
\(A(v)\le v^2d_*\), \(K>2\), and the negligible artificial
endpoint term. For that last point,
\(d_*\le Cd s_0^{-m}\) implies
\(s_0\le C\exp(-(T-a_0)/m)\ll\rho\).
At the outer endpoint the artificial term dominates, and the exact
formula is
\[
 a(\rho)=T+(2m+4)\log\rho=T-o(T)>T/2.
\]
Before the first crossing of \(T/2\), the nonendpoint terms have
\(I(s)/s^2\le d_*\), and the endpoint term has
\(I_{\mathrm{end}}(s)/s^2\le d\rho^{-2m-4}=o(d_*)\).
This proves the remaining assertions about the levels.

On any compact logarithmic scale interval, \(A\) is smooth and has
bounded derivative. The supremum of
\(v^{-K}A(v)\) over \([s,\rho]\), with the positive endpoint floor,
is therefore locally Lipschitz. Its derivative vanishes unless its
maximum is realized at the moving endpoint \(v=s\). On the equality
set, the derivatives of \(I\) and \(A\) agree almost everywhere.
This proves the reset rule. Gaussian differentiation, the dilation
bound, and \Cref{lt:gaussian} give
\(|\dot A|+|\dot Z|+|\ddot Z|\le CI\). The constant \(d\) term
is negligible, since \(s^2\eta d/I=\eta d e^a\le e^{-T/2}\).
Together with the reset rule these prove \eqref{ex:derivatives}
and the assertion about \(\dot Z/I\).

Finally all polynomial dilations in \(a\) used here lie well inside
the chart. Indeed
\(I(s)\le Cd s^{2-m}+I_{\mathrm{end}}(s)\), whereas
\(I(s)\ge s^2e^{-T/2}\) on the path. The endpoint term cannot
account for this inequality. Thus \(s\le Ce^{-T/(2m)}\), up to
an immaterial fixed weakening of the exponent. This proves the
claimed interior room and justifies the uniform Gaussian cutoffs.
\end{proof}

\subsection{Height multipliers and initialization}

For \(Q>1\), choose an increasing smooth function \(b\) with
\[
 b(x)\ge (\log x)^{201}
       +3F_{Q-1}^{-1}((\log x)^3),\qquad b(x)=x^{o(1)},
\]
using this displayed sum up to fixed factors. Write
\(S_1(x)=\exp(Cb(x))\). Its constant is large enough to dominate
all polynomial and dilation losses in \Cref{ct:separated}, including
applications at \(x-O(\log x)\). A further fixed power is denoted by
\(S_h(x)\). When \(Q=1\), \(S_h\) denotes only a sufficiently large
fixed constant; no function \(b\), separated estimate, or multiple
well construction is used.

All configuration lists and joint moment or support witnesses below
use the same fixed exceptional proportion \(\delta_{\rm maj}\) when
\(Q>1\); for \(Q=1\), fix one sufficiently small proportion as well.
These proportions are chosen independently of \(T\) and \(a\), and
their corresponding small-density thresholds in \Cref{lt:config}
are then fixed. The densities used below are bounded by fixed
polynomials times \(e^{-a_0}\) or \(e^{-a}\), with
\(a\ge a_0-O(1)\); hence they lie below those thresholds for large \(T\).

Here is the precise multiplier convention. For a finite collection
\(\mathcal W=\{b_1,\ldots,b_N\}\) of distinct wells let
\(d_{\mathcal W}(z)=\dist(z,\mathcal W)\). There is a smooth function
\(\Gamma_0\), comparable to \(d_{\mathcal W}^2\), equal to
\(|z-b_i|^2\) within a fixed small fraction of the nearest-neighbor
gap at each \(b_i\), and satisfying
\[
 |D\Gamma_0|^2\le C\Gamma_0,\qquad |D^2\Gamma_0|\le C.
\]
Construct it by the Whitney partition method \cite{Whitney1934}
on the complement of the wells:
on a Whitney ball of radius comparable to distance from the wells,
the squared distance has size \(O(r^2)\), first variation \(O(r)\),
and the partition derivatives have sizes \(O(r^{-1})\) and
\(O(r^{-2})\). Patch to the exact quadratics on their disjoint near
neighborhoods. This gives the asserted uniform derivative bounds.
For one well take its exact square.

Choose a fixed smooth nondecreasing clipping function \(\chi\),
equal to its argument near zero and constant for large arguments,
with \(\chi(t)\) comparable to \(\min(t,1)\), and set
\begin{equation}
 \Gamma_s(z)=s^2\chi(\Gamma_0(z)/s^2),\qquad
 U(s)=\int G_s(y-x)\Gamma_s(z)\,\dd\mu(y,z).
 \label{ex:moment}
\end{equation}
We may arrange a smooth vector multiplier \(V_b\), equal to
\(z-b_i\) in smaller near neighborhoods and vanishing outside
the exact-quadratic neighborhoods and at distances comparable to
\(s\), such that
\begin{equation}
 |V_b|^2\le\Gamma_s,\qquad |DV_b|\le C,\qquad
 |D\Gamma_s|^2\le C\Gamma_s,\qquad |D^2\Gamma_s|\le C.
 \label{ex:multiplier-bounds}
\end{equation}
Use disjoint radial cutoffs of \(z-b_i\); their differentiated
cutoffs multiply a displacement of the same size as their radius.

The transition set includes every point outside the smaller pure
neighborhoods where either \(D^2\Gamma_s\ne2\operatorname{Id}\)
or \(DV_b\ne\operatorname{Id}\). In particular it includes the
cutoff annuli of \(V_b\) even if \(\Gamma_s\) is still exactly
quadratic there. Its Gaussian tilt is denoted by \(D_{\mathrm{tr}}\).
We account separately for Gaussian tails on which clipping occurs.
In the one-well case all derivative-error regions lie in those tails,
and \(D_{\mathrm{tr}}\) is omitted.

To state the initialization, write $N$ for the number of wells and,
for a small positive parameter $k$, put $q=-k/a^2$. The comparison to
be maintained is
\begin{equation}
 kU+\alpha Z\le I,\qquad
 \alpha=
 \begin{cases}
 1+C_1ka/\log a,&N>1,\\
 1+(1-4k)q-2k,&N=1,
 \end{cases}
 \label{ex:comparison}
\end{equation}
where $C_1$ is a sufficiently large fixed constant. The next lemma
chooses the wells and $k$ so that this inequality holds at $s_0$.

\begin{lemma}[Initialization at a threshold or a side]
\label{ex:initialization}
At \(s_0\) one can choose at most \(Q\) wells and a small parameter
\(k_0>0\) so that the barrier of \Cref{ex:barrier} holds initially.
If more than one well is used, one may arrange, for any fixed \(P\),
\[
 Z(s_0)\le a_0^{-P}I(s_0),\qquad U(s_0)\le CI(s_0),
\]
before any immediate mergers, with \(k_0\) a small fixed constant.
Alternatively there is one common well, with
\[
 U(s_0)\le CS_h(a_0)^2 I(s_0),\qquad
 I(s_0)-Z(s_0)\ge cS_h(a_0)^{-2}I(s_0),
 \qquad k_0=cS_h(a_0)^{-6}.
\]
For \(Q=1\), \(k_0\) is a small fixed constant.
\end{lemma}

\begin{proof}
The bound on \(s_0\) in the envelope proof gives
\(s_0\operatorname{poly}(a_0)\ll\rho\). Threshold control on fixed
enlargements therefore selects \(Q\) by \Cref{lt:scalar}. Choose a
joint majority-and-moment witness
\(\mathcal A_0=\{a_1,\ldots,a_Q\}\) from \Cref{lt:config} at scale
\(s_0\), using \(B_{r_{\rm maj}s_0}(x)\) and the fixed majority
proportion when \(Q>1\), and a fixed concentric inner ball when
\(Q=1\). Use its distinct entries as the initial wells. Put
\(R_0=C_b\sqrt{a_0}\), with a sufficiently large fixed \(C_b>2\).
All fixed enlargements
needed on \(|u|\le R_0\), where \(u=(y-x)/s_0\), lie below \(\rho\).
Threshold control gives normalized density at most \(Cd_*\) on these
dilations. The dilation part of \Cref{lt:weightedballs}, applied only
on this body, and
\(\Gamma_{s_0}\le C\dist(z,\mathcal A_0)^2\) give
\[
 \int_{|u|\le R_0}G_{s_0}\Gamma_{s_0}\,\dd\mu
 \le Cs_0^2d_*.
\]
The root mass identity and threshold control at \(\sqrt2s_0<\rho\)
give
\[
 (\pi_\#\mu)_{\sqrt2s_0}
 \le C+\tau_{\sqrt2s_0}
       +Cs_0^{-m}e^{-c/s_0^2}\le C.
\]
The last term accounts for the fixed exterior-chart and periodic
tails. Since \(\Gamma_{s_0}\le Cs_0^2\), Gaussian comparison yields
\[
 \int_{|u|>R_0}G_{s_0}\Gamma_{s_0}\,\dd\mu
 \le Cs_0^2e^{-R_0^2/4}(\pi_\#\mu)_{\sqrt2s_0}
 \le Cs_0^2e^{-C_b^2a_0/4}=o(s_0^2d_*).
\]
For the periodized kernel the split and comparison are made termwise
in its lifts. Consequently
\(U(s_0)\le Cs_0^2d_*=O(I(s_0))\).
If the contact itself is a threshold equality, then
\(Z(s_0)=0\), which gives the first initialization, with one well
when \(Q=1\).

Assume now that \(\underline s<s_0\). For
\(\underline s\le v\le s_0\), put \(R_y=|y-x|/v\) and consider the
Gaussian work cell \(B_r(y)\) with \(r=cv/(1+R_y)\), where
\(R_y\le C\sqrt{a_0}\). For any fixed required enlargement \(B_{Lr}\),
the Gaussian at \(C_L(1+R_y)v\) is bounded below by a fixed multiple
of \(((1+R_y)v)^{-m}\) on that enlargement. This scale lies in
\([\underline s,\rho]\), so the threshold bound gives, in particular,
\begin{equation}
 \begin{split}
 r^{-m}\tau(B_{4r}(y))&\le d_c:=C(1+R_y)^{2m}d_*,\\
 T_c:=-\log d_c&=a_0-2m\log(1+R_y)-O(1),\\
 r\sqrt{d_c}&\le Cv(1+R_y)^{m-1}\sqrt{d_*}.
 \end{split}
 \label{ex:threshold-cells}
\end{equation}
The same estimate holds for \(B_{Lr}\) after changing its fixed
constant. Since \(d_c=o(1)\), it supplies small \(D+\Lambda\)
energy and bounded normalized projected mass. Its \(|M|\) part puts
the normalized smooth mass average within \(Cd_c<1/4\) of \(Q\),
so \Cref{lt:scalar} selects that integer on every cell. The cell
radius may be below \(\underline s\); the threshold is used only at
the larger Gaussian scale.

When \(Q>1\), scale \(B_r(y)\) to one and choose a length-\(Q\)
majority list \(\mathcal A_{v,y}\) on \(B_{r_{\rm maj}}\) with local
upper bound \(r_{\rm maj}^{-m}d_c\) and the fixed
\(\delta_{\rm maj},C_{\rm maj}\), as required by
\Cref{ct:separated}. Multiply its heights by \(r\) when comparing
them in the fixed coordinates of this section. The overlap and
matching rules \eqref{ct:majority-overlap}--\eqref{ct:majority-matching}
match these lists to \(\mathcal A_0\). Along concentric dyadic scales
from \(s_0\) to \(v\), the increments \(Cr'\sqrt{d_*}\) at radius
\(r'\) sum to \(Cs_0\sqrt{d_*}\). At the fixed scale \(v\), enlarge
to radius \(C(1+R_y)v\), move the center in the prescribed steps,
and shrink to \(r\). All fixed enlargements in this part lie inside
\(B_{C'(1+R_y)v}(x)\); at an intermediate radius \(r'\) their
normalized densities are at most
\(C((1+R_y)v/r')^m d_*\). The matching increments therefore add
only \(C(1+R_y)^C v\sqrt{d_*}\). The total motion is bounded by a
fixed polynomial in \(a_0\) times \(s_0\sqrt{d_*}\), uniformly even
as the bottom tends to zero.

Suppose \(\mathcal A_0\) has at least two nonempty groups separated
by more than \(S_h(a_0)s_0\sqrt{d_*}\). Group the \(Q\) entries of
each \(\mathcal A_{v,y}\) according to the corresponding separated
groups of the initial wells. Matching preserves their nonemptiness,
and their gaps exceed \(10S(T_c)r\sqrt{d_c}\): the fixed power
\(S_h(a_0)\) absorbs \(S(T_c)\), the displayed cell factor, and the
matching loss, since \(T_c=a_0-O(\log a_0)\). These length-\(Q\)
lists are the inputs to \Cref{ct:separated}. The cell tests obtained
by partitioning the nonnegative Gaussian are admissible after
normalization in the cell coordinates. Bounded overlap and a
sufficiently large input exponent absorb their fixed polynomial
losses. In the curved case the physical cell radius is \(lr\),
and \(T_c=o(T)\), so the original physical smallness implies the
cell hypothesis. The threshold and root tail bounds from the
envelope proof apply with controlling density \(d_*\), uniformly
down to \(\underline s\). Under nonnegative testing,
\(f_*\le L-C_Q\Lambda\), and the far term in
\Cref{ct:separated} is nonnegative. The resulting bound, integrated
in scale, gives
\[
 Z(s_0)\le Ca_0^{-P-1}d_*
           \int_{\underline s}^{s_0}v^2\,\dd\log v
 \le a_0^{-P}I(s_0).
\]
Entries too close to remain separate are merged immediately by the
rule below. Their moment before merging is \(O(I(s_0))\).

If there are no such separated groups, a chain through at most
\(Q\) entries of \(\mathcal A_0\) shows that its diameter is at most
\(CQS_h(a_0)s_0\sqrt{d_*}\). Use one common well.
Pointwise, the new clipped square is bounded by a fixed multiple of
the old one plus the squared diameter. The preceding \(U\)-bound and
bounded Gaussian mass therefore give the stated estimate for \(U\).
It remains to check the surplus of \(I\) over \(Z\). If
\(A(s_0)\) is a sufficiently small fraction of \(I(s_0)\), this is
immediate from \eqref{ex:bank}. Otherwise
\(A(s_0)\ge cI(s_0)\). When \(\underline s\ge s_0/2\), the
inequality \(\dot A\le2A+s^2R_s+\text{error}\) gives
\[
 A(\underline s)+\int_{\underline s}^{s_0}v^2R_v\,\dd\log v
 \ge cI(s_0).
\]
These terms are retained in \(A(s_0)-Z(s_0)\), after absorbing the
error into the positive margin.

When \(\underline s<s_0/2\), Gaussian differentiation and
\eqref{ex:threshold-control} give
\(|\dot A(v)|\le Cv^2\tau_{2v}\le Cv^2d_*\) on \([s_0/2,s_0]\).
Since \(A(s_0)\ge cI(s_0)\asymp s_0^2d_*\), for some fixed
\(0<\vartheta_0<\log2\) one has \(A(v)\ge cI(s_0)\) throughout
\([e^{-\vartheta_0}s_0,s_0]\). Apply \Cref{lt:weightedballs} directly
on this earlier interval, using the scalar hypotheses verified in
\eqref{ex:threshold-cells}. The tail implication in the envelope
proof applies with controlling density \(d_*\): here
\(v<\rho/T^N\) for every fixed \(N\), and \(d_*\ge e^{-T/2}\).
It gives an \(o(d_*)\) remainder uniformly in the bottom. Thus
\(|M|_v\le C(D_v+\Lambda_v)+o(d_*)\); because
\(M_v=A(v)/v^2\ge cd_*\), this yields \(D_v+\Lambda_v\ge cd_*\).

The surplus uses further information. Where \(\Lambda_v\ge cd_*\),
the retained term \(C_Qv^2\Lambda_v\) in \(J(v)\) pays it directly.
On the complementary portion \(D_v\ge cd_*\). Let \(V_v\) be
displacement from the common well, smoothly clipped at distance
comparable to \(v\). Apply the moment part of \Cref{lt:weightedballs}
to joint configuration witnesses on the admissible dilations and
match their centers to the common well. Gaussian summation absorbs
their polynomial losses and gives
\[
 \int G_v|V_v|^2\,\dd\mu
 \le Cv^2S_h(a_0)^2d_*.
\]
The support estimates for those chosen witnesses put all support
heights on \(|y-x|\le C_b\sqrt{a_0}v\) within
\[
 v\bigl[\operatorname{poly}(a_0)e^{-\beta a_0}
              +S_h(a_0)e^{-a_0/2}\bigr]=o(v)
\]
of the common well. Clipping is therefore absent on this body.
For the clipped tail, \(\pi_\#\mu=Q\vol_S+M\) controls the bare
mass as well as the defect. Since \(\log S_h(a_0)=o(a_0)\), taking
the fixed \(C_b\) sufficiently large makes the Gaussian mass tail
\(o(d_*S_h(a_0)^{-2})\). The defect and chart tails obey the
preceding uniform bounds.

Choose one fixed root-chart cutoff \(0\le\chi\le1\), equal to one on
the fixed working region, whose transition is at distance at least
\(h\) from that region and has width at least \(ch\), decreasing \(h\)
to the available width. Then
\(\|\nabla\chi\|_\infty\le C(1+h^{-1})\); the first-variation product
rule differentiates \(\chi\) once, while all scale derivatives leave
it fixed.
With \(u=(y-x)/v\), normal stationarity in \Cref{lt:heighttests}
now gives
\[
 D_v\le \frac1v
       \left|\int G_v V_v\cdot P_{vh}u\,\dd\mu\right|+o(d_*).
\]
The factor \(v^{-1}\) is the Gaussian derivative. Retaining
\(|V_v|\) in the curved normal error and using its moment bound
leaves relative error
\(e^{-\kappa T}\operatorname{poly}(a_0,S_h(a_0))=o(1)\);
the normal identity introduces no further inverse power of \(v\).
Thus the flux has size at least \(cvd_*\) on the \(D\)-portion.
Cauchy--Schwarz and \(|P_{vh}u|^2\le B[u,u]\) yield
\[
 c v^2d_*^2
 \le \left(\int G_v|V_v|^2\,\dd\mu\right)R_v
 \le Cv^2S_h(a_0)^2d_*R_v,
 \qquad
 v^2R_v\ge cS_h(a_0)^{-2}v^2d_*.
\]
After the margin absorbs the errors, \(\dot A-\dot Z\) retains
\(v^2R_v+C_Qv^2\Lambda_v\); at all other scales it is nonnegative
by \eqref{ex:bank}. Integration over the fixed logarithmic interval
therefore proves the surplus claimed in the statement. Since the
one-well coefficient \(\alpha\)
below is less than one, \(k_0=cS_h(a_0)^{-6}\) pays for the moment
while leaving a strict initial inequality. The same proof for
\(Q=1\) uses the Gaussian-summed moment with fixed \(S_h\) and no
separation step.
\end{proof}

\subsection{The barrier and its finite mergers}

Normal stationarity and square completion bound \(s^2R_s\) from below
in terms of \(D_s\) and \(U\), supplying the term that offsets growth
of \(kU\). With several wells, \(\alpha\) makes the coefficient of
\(\dot Z\) in the derivative of \(I-kU-\alpha Z\) negative, so that
the separated far-tilt gain absorbs the transition-tilt
losses. The merger rule preserves this comparison with controlled losses,
leaving the traversal required by \Cref{ex:bands}.

\begin{proposition}[Low-frequency barrier]
\label{ex:barrier}
Along every path in \Cref{ex:envelope}, the wells can be chosen
piecewise constantly, with at most \(Q-1\) mergers. There is a
positive piecewise smooth parameter \(k\), bounded above by a small
fixed constant and satisfying between mergers
\begin{equation}
 \frac{k_a}{k}=q=-\frac{k}{a^2},
 \label{ex:k-evolution}
\end{equation}
such that the comparison \eqref{ex:comparison} holds.
Here \(C_1\) is fixed and sufficiently large. Uniformly along the
path, \(\log(1/k)=o(a_0)\) as \(T\to\infty\).
Before any hypothetical first failure of the barrier, the following
estimates hold, with \(\mathcal E\le a^{-P_0/2}\):
\begin{align}
 \dot U&\le2s^2D_s+Cs^2D_{\mathrm{tr}}+\mathcal E I,
 \label{ex:moment-derivative}\\
 s^2R_s&\ge4k(s^2D_s-Cs^2D_{\mathrm{tr}})
             -4k^2U-\mathcal E kI.
 \label{ex:radial-bound}
\end{align}
With more than one well, in addition,
\begin{equation}
 \dot Z+c_2\frac{(\log a)^2}{a}s^2D_{\mathrm{tr}}
       \le a^{-P_0}I.
 \label{ex:transition}
\end{equation}
The transition terms are absent with one well. By increasing the
fixed exponents in the construction, \(\mathcal E\) can be smaller
than any prescribed inverse power of \(a\).
\end{proposition}

\begin{proof}
We specify the mergers first. At stage \(j\), merge a closest pair
as soon as its gap \(g\) satisfies
\begin{equation}
 g\le s\min\{C_j e^{-\gamma a},\,
                   S_1(a)e^{-a/2}k^{-1}\}.
 \label{ex:merge-rule}
\end{equation}
Retain one of the two centers. Take \(\gamma>0\) sufficiently small,
\(C_{j+1}\) a sufficiently large fixed multiple of \(QC_j\), and
replace \(k\) by
\begin{equation}
 k_{\mathrm{new}}=c k^6S_1(a)^{-6}.
 \label{ex:merge-k}
\end{equation}
Perform further mergers immediately if required. For \(Q=1\)
there is no merger rule. A common-well initialization when \(Q>1\)
may be counted as an initial merger; it uses no more than the
available \(Q-1\) losses.

The solution of \eqref{ex:k-evolution} on a stage starting at
level \(a_j\) is
\[
 k(a)^{-1}=k(a_j)^{-1}+a_j^{-1}-a^{-1}.
\]
Retracing levels therefore causes no accumulated loss. Since
\(a\ge a_0-O(1)\), \(k\) changes by at most a fixed factor on a
stage. There are at most \(Q-1\) mergers, and their levels are at
most \(T/2\). Thus
\[
 \log(1/k)\le C_Q\{1+b(T)+b(a_0)\}=o(a_0),
\]
where \Cref{ex:bands} verifies the final comparison. In particular
\(e^{-a}/k\) is exponentially small uniformly along the path.

We next verify that the wells continue to represent the configuration.
At a scale \(s\), let \(\mathcal A_s\) be a length-\(Q\) majority list
on a fixed inner ball, using the fixed convention above when \(Q>1\).
For support control, choose joint majority-and-support witnesses from
\Cref{lt:config} on the balls under consideration and match their
centers to these lists. The support and matching bounds, together with
\eqref{ex:dilations}, bound configuration changes on a ball of scale
\(v\), including polynomial enlargements, by
\(Cv e^{-\beta a(v)}\) for some fixed \(\beta>0\).
Decrease \(\beta\) if necessary so that
\(\beta|\dot a|\le1/2\). Summing on a geometric sequence from
\(s_0\) to \(s\) gives \(Cs e^{-\beta a(s)}\). The same estimate
tracks each initial configuration entry. The additional displacement
caused by previous mergers is bounded by the sum of their gaps.
Choose \(\gamma\ll\beta\) and \(\gamma|\dot a|\le1/2\). Then
\(s e^{-\gamma a(s)}\) increases outwards, and the successive
enlargements of \(C_j\) make all earlier capped gaps a small fixed
fraction of the current cap. These facts give, in particular,
\begin{equation}
 \dist(z,\mathcal W)\le Cs e^{-\gamma' a}
 \quad\text{on }|y-x|/s\le C\sqrt a
 \label{ex:rough-height}
\end{equation}
for some fixed \(\gamma'>0\). This estimate does not use the
barrier. It also bounds the distance from retained wells to the
corresponding tracked configurations. The common-well initialization
has the same property, since \(\log S_h(a_0)=o(a_0)\).

While the barrier holds, each entry of \(\mathcal A_s\) is matched on
a fixed majority of the inner ball to a projection entry within
\(Cs e^{-a/2}\). The area formula and the uniform upper bound for
\(J_\pi\) give that entry a fixed amount of Gaussian mass. The bound
\(|D\Gamma_s|^2\le C\Gamma_s\) and the numerical barrier therefore give
\[
 \max_{\xi\in\mathcal A_s}\Gamma_s(\xi)
 \le C U+Cs^2e^{-a}
 \le Cs^2e^{-a}/k.
\]
Since \(\Gamma_s\) is comparable to
\(\min\{\dist(z,\mathcal W)^2,s^2\}\) and \(e^{-a}/k\to0\), every
entry lies within \(Cs\sqrt{e^{-a}/k}\) of a well. This distance is a small
fraction of \(S_1(a)se^{-a/2}k^{-1}\). It is also a small fraction
of the cap, using \(\log(1/k)=o(a_0)\); alternatively
\eqref{ex:rough-height} gives that conclusion directly. When no
merger is triggered, the resulting neighborhoods of the wells are
disjoint. The group represented by a well is the sublist of
\(\mathcal A_s\) in its neighborhood. When \(Q>1\), its integer count
cannot change: on two sufficiently nearby scales the majority balls
obey \eqref{ct:majority-overlap}, their lists match by
\eqref{ct:majority-matching}, and the matching error is much smaller
than the gap. For \(Q=1\), the sole length-one group remains represented.
At initialization each well is represented. At a merger
the surviving neighborhood contains the entries of both merged wells
after any further immediate mergers. This proves representation at
every stage, including a first barrier equality.

With several wells, \(Q>1\), so we may apply \Cref{ct:separated}.
For a Gaussian cell with \(R_y=|y-x|/s\le C\sqrt a\) and
\(r=cs/(1+R_y)\), the dilation bound at \(C(1+R_y)s\) gives
\begin{equation}
 \begin{split}
 r^{-m}\tau(B_{4r}(y))
   &\le d_c:=C(1+R_y)^{2m+K-2}e^{-a},\\
 T_c:=-\log d_c
   &=a-(2m+K-2)\log(1+R_y)-O(1),\\
 r\sqrt{d_c}
   &\le Cs(1+R_y)^{m+K/2-2}e^{-a/2}.
 \end{split}
 \label{ex:transition-cells}
\end{equation}
The cell and its fixed enlargements lie in the chart, and
\(T_c=a-O(\log a)\to\infty\). As in \eqref{ex:threshold-cells},
the \(|M|\) bound selects \(Q\). In the cell coordinates choose a
length-\(Q\) majority list \(\mathcal A_{s,y}\) on
\(B_{r_{\rm maj}}\), with local upper bound \(r_{\rm maj}^{-m}d_c\)
and the fixed \(\delta_{\rm maj},C_{\rm maj}\), and multiply its
heights by \(r\) for comparison in the section coordinates. The
off-center overlap chain from initialization now uses
\eqref{ex:dilations}; matching to \(\mathcal A_s\) costs at most a
fixed polynomial in \(a\) times \(s e^{-a/2}\).

Group the \(Q\) entries of \(\mathcal A_{s,y}\) by the represented
wells to which their matches in \(\mathcal A_s\) belong. The wells
index these nonempty groups; the input list still has length \(Q\)
after mergers. Before a merger, if the second threshold in
\eqref{ex:merge-rule} is smaller, the gap is larger than
\(S_1(a)s e^{-a/2}k^{-1}\). This dominates the entry-to-well and
cell matching errors and \(10S(T_c)r\sqrt{d_c}\), by the choice of
\(S_1\). If the cap is smaller, the ratio of its gap to the last
quantity tends to infinity, since it is bounded below by
\(\exp((1/2-\gamma)a-O(b(a))-O(\log a))\). Thus the groups meet
the separation premise in both regimes. Away from the clipping tails,
every point of the transition set is a fixed fraction of a gap from
the represented groups, and remains in the cell's far region after
matching. This includes every extra annulus in \(D_{\mathrm{tr}}\).

Apply \Cref{ct:separated} to the nonnegative normalized cell weights.
Their polynomial Lipschitz costs are absorbed by taking its exponent
larger than \(P_0\), while
\(c(T_c)\asymp(\log a)^2/a\). Bounded overlap and the Gaussian tails
give total remainder \(O(a^{-P_0}I/s^2)\). The curved hypothesis is
inherited because the physical cell radius is \(lr\) and
\(T_c\le T/2+O(1)\). This proves \eqref{ex:transition}.

We prove the two differential inequalities explicitly. In flat
coordinates the horizontal test with multiplier \(\Gamma_s\) gives
\[
 \dot U=\int G_s\{\Gamma_s(B[u,u]-e)
             +sD\Gamma_s\cdot P_{vh}u+\dot\Gamma_s\}\,\dd\mu,
 \qquad u=(y-x)/s.
\]
Normal stationarity with \(G_sD\Gamma_s\) converts the cross term
to \(s^2\int G_sD^2\Gamma_s:P_{vv}\,\dd\mu\).
On the pure regions \(D^2\Gamma_s=2\operatorname{Id}\), and
elsewhere its norm is bounded. Dropping \(-\Gamma_se\), and using
\eqref{ex:rough-height} to bound the \(\Gamma_s B[u,u]\) term,
gives \eqref{ex:moment-derivative}. The explicit derivative
\(\dot\Gamma_s=2s^2[\chi(t)-t\chi'(t)]\),
\(t=\Gamma_0/s^2\), is supported on the clipping region. That region
does not meet the Gaussian body by \eqref{ex:rough-height}.

For the second inequality use \(|P_{vh}u|^2\le B[u,u]\) and
complete the square:
\[
 s^2R_s\ge
 4ks\int G_sV_b\cdot P_{vh}u\,\dd\mu-4k^2U.
\]
Normal stationarity gives
\[
 s\int G_sV_b\cdot P_{vh}u\,\dd\mu
   =s^2\int G_sDV_b:P_{vv}\,\dd\mu
   \ge s^2D_s-Cs^2D_{\mathrm{tr}}
\]
up to the already specified tails. This proves
\eqref{ex:radial-bound} in flat coordinates.

For completeness the tails here also include terms involving mass
without a tilt factor. Since \(\Gamma_s\le Cs^2\), the mass tail
beyond \(C\sqrt a\) is bounded by
\(Cs^2e^{-cC^2a}\), plus the dilation-controlled defect. Taking
\(C\) large makes this smaller than any required multiple of
\(a^{-P}I\). Terms with tilt use \eqref{ex:dilations} directly.
No inverse power of \(k\) is required for these flat errors, because
the derivative bounds in \eqref{ex:multiplier-bounds} are uniform.
For the curved errors, keep the original coordinates scaled by the
physical radius \(l\), and write \(\langle f\rangle_s=\int G_sf\,\dd\mu\)
and \(\gamma_{\rm geom}=l+\sup|z_{\mathrm{phys}}|\). The dilation bound gives
\(\langle(1+|u|^2)(e+\lambda)\rangle_s\le P(a)I/s^2\),
where \(P\) is a fixed polynomial and
\(\lambda=|z_{\mathrm{phys}}|^2+|H_S|^2+|\nabla H_S|^2\).
For either \(v=V_b\) or \(v=D\Gamma_s\),
\(\langle|v|^2\rangle_s\le CU\le CI/k\) and \(|Dv|\le C\).
After multiplying the normal identity by \(s^2\), the errors in
\Cref{lt:heighttests} satisfy
\begin{align*}
 s^2l\langle|z_{\mathrm{phys}}|\sqrt e\rangle_s
 &\le s^2l\sqrt{\Lambda_sD_s}\le CP(a)lI,\\
 s^2l\langle|v|(|H_S|+|z_{\mathrm{phys}}|+e+\sqrt e)\rangle_s
 &\le Cs^2l\sqrt U\sqrt{\Lambda_s+D_s}
       \le CP(a)ls k^{-1/2}I,\\
 \gamma_{\rm geom} s\langle|u||v|\sqrt e\rangle_s
 &\le\gamma_{\rm geom} s\sqrt U\sqrt{\langle|u|^2e\rangle_s}
       \le CP(a)\gamma_{\rm geom} k^{-1/2}I.
\end{align*}
Here \(e\le m\) controls the term \(|v|e\); the last line is
the error from differentiating the Gaussian. In the horizontal test,
\(\Gamma_s\le Cs^2\) bounds the error carrying
\(\Gamma_s(e+\lambda)\) by \(C\gamma_{\rm geom} P(a)I\), while
\[
 \gamma_{\rm geom} s\langle|u||D\Gamma_s|(\sqrt e+\sqrt\lambda)\rangle_s
 \le C\gamma_{\rm geom} s\sqrt U
             \sqrt{\langle|u|^2(e+\lambda)\rangle_s}
 \le C\gamma_{\rm geom} P(a)k^{-1/2}I.
\]
Coordinate and base-volume comparison may also produce a term
\(C\gamma_{\rm geom}\langle(1+|u|^2)\Gamma_s\rangle_s\).
Split at \(C\sqrt a\); its body is bounded by
\(C\gamma_{\rm geom} P(a)U\), and its clipped tail has the already proved
negligible bound. Thus it costs at most
\(C\gamma_{\rm geom} P(a)k^{-1}I\). These estimates retain physical heights
and curvature inside \(\Lambda\) until after Cauchy--Schwarz;
the displayed factors of \(s\) cancel the Gaussian inverse-scale
factors. They give total error
\(Ce^{-\kappa T}\operatorname{poly}(a,k^{-1})I\) in the moment
inequality and this bound times \(k\) after square completion in
the radial inequality. Since
\(\log(1/k)=o(T)\), these errors are smaller than the stated
\(\mathcal E I\) and \(\mathcal E kI\). Chart-cutoff errors are
still smaller because \(s\le e^{-cT}\). This proves both
differential inequalities with their asserted uniform errors.

It remains to prove the comparison, including its signs. Work at
an equality \(kU+\alpha Z=I\) on a fixed stage and put
\[
 \delta_\alpha=\alpha_a-q\alpha,\qquad
 h_1=q+\delta_\alpha Z/I.
\]
Differentiating \(I-kU-\alpha Z\), and using
\(\dot a=2-\dot I/I\), gives
\begin{equation}
 (I-kU-\alpha Z)^{\boldsymbol\cdot}
 =(1+h_1)\dot I-2h_1I-k\dot U-\alpha\dot Z.
 \label{ex:barrier-derivative}
\end{equation}
At a reset write \(\mathcal D=s^2D_s\),
\(\mathcal R=s^2R_s\). Apart from the absorbed geometric error,
\[
 \dot A=2I-\mathcal D+\mathcal R,\qquad
 \mathcal D=2I-\dot Z-J.
\]
Insert \eqref{ex:moment-derivative} and
\eqref{ex:radial-bound} into \eqref{ex:barrier-derivative} and
use \(kU=I-\alpha Z\). The resulting lower bound is
\begin{align}
 &(4k-2)h_1 I+4k(1+h_1)\alpha Z
    +c_Z\dot Z-Cks^2D_{\mathrm{tr}}-O(\mathcal E kI)
       +c_JJ,\label{ex:reset-algebra}\\
 &c_Z=1+h_1-2k-4kh_1-\alpha,\qquad
 c_J=2k-(4k-1)(1+h_1).\notag
\end{align}
In both cases \(h_1=o(1)\), so \(c_J>0\); discarding the
margin therefore has the indicated favorable sign. The first-variation
error in \(\dot A\), which initially need not contain \(k\), is
absorbed by a fixed fraction of this positive margin. Equivalently,
its remaining contribution divided by \(kI\) can be included in
\(\mathcal E\), since \(\log(1/k)=o(T)\).

With one well,
\[
 \delta_\alpha=\frac{k}{a^2}
       +\frac{2k(1-4k)}{a^3}-\frac{4k^3}{a^4}=O(k/a^2),\qquad
 c_Z=(1-4k)\delta_\alpha Z/I.
\]
Thus the \(\dot Z\) coefficient cancels exactly at \(Z=0\).
The estimate \(|\dot Z|\le CI\) bounds its remaining contribution
by \(O(k/a^2)Z\). The term \((4k-2)qI\) is at least
\(ckI/a^2\); the coefficient of \(Z\) left in the first two
terms of \eqref{ex:reset-algebra} is at least \(ck\).
Hence the reset derivative is strictly positive for large \(a\).

With several wells there is the useful exact cancellation
\[
 \delta_\alpha=\frac{k}{a^2}
       +C_1k\frac{\log a-1}{(\log a)^2}.
\]
In particular \(\delta_\alpha=o(k)\),
\(|h_1|\le C/a\), and
\[
 -C C_1ka/\log a\le c_Z\le-c C_1ka/\log a.
\]
Multiply \eqref{ex:transition} by this negative coefficient.
It supplies a favorable multiple of
\(C_1k\log a\,s^2D_{\mathrm{tr}}\), which absorbs the
transition loss in \eqref{ex:reset-algebra}. Its adverse error is
\(O(k a^{1-P_0}/\log a)I=o(kI/a^2)\). The remaining \(Z\)
terms are positive, since their possible loss is
\(O(\delta_\alpha)Z=o(k)Z\). Again
\((4k-2)qI\ge ckI/a^2\) supplies strict positivity.

Off the reset set \(\dot I=KI\), and
\(\mathcal D\le2I-\dot Z\). The coarser consequence
\(\dot U\le C\mathcal D+\mathcal E I\) gives
\[
 (I-kU-\alpha Z)^{\boldsymbol\cdot}
 \ge\{(1+h_1)K-2h_1-2Ck\}I
          +(Ck-\alpha)\dot Z-O(\mathcal E kI).
\]
The coefficient \(Ck-\alpha\) is negative. A negative \(\dot Z\)
therefore helps. For one well, \(\dot Z\le2I\); for several wells,
\eqref{ex:transition} gives \(\dot Z\le a^{-P_0}I\), and
\(\alpha a^{-P_0}=o(1)\). Taking \(K\) large and the initial
upper bound for \(k\) small makes the derivative strictly positive.
All the strict inequalities hold in a neighborhood of equality on
each compact stage. The absolutely continuous comparison principle
therefore applies despite the reset set having only almost-everywhere
derivatives. Initialization is supplied by
\Cref{ex:initialization}; in its multiple-well case choose
\(P>3\), so that the initially large coefficient \(\alpha\)
still leaves strict room.

Finally consider a merger. Pointwise distance comparison and bounded
Gaussian mass imply
\[
 U_{\mathrm{new}}\le C(U+g^2),\qquad
 g^2\le S_1(a)^2k^{-2}I.
\]
Moreover \(\alpha_{\mathrm{new}}\le(1-ck)\alpha\), whether
several wells remain or the last two merge: in the first case use
the formula linear in \(k\). In the second case write
\(\alpha=1+x\), where \(x=C_1ka/\log a\). The fixed upper bound
for \(k\) and large \(a\) give \(x/(1+x)\ge ck\), so
\(\alpha_{\mathrm{new}}<1\le(1-ck)\alpha\).
Set \(\theta=Ck_{\mathrm{new}}/k\) and
\(\lambda=Ck_{\mathrm{new}}S_1^2k^{-2}\). By
\eqref{ex:merge-k}, \(\theta\ll1\) and \(\lambda\ll k\).
The old comparison then gives
\[
 k_{\mathrm{new}}U_{\mathrm{new}}+\alpha_{\mathrm{new}}Z
 \le(\theta+\lambda)I+
        (\alpha_{\mathrm{new}}-\theta\alpha)Z\le I.
\]
For the last inequality take the maximum over
\(0\le Z\le I/\alpha\); both endpoint values are below \(I\).
Thus every merger preserves the barrier, completing its proof.
\end{proof}

\subsection{Uniform bands of levels}

\begin{lemma}[Choice of bands]
\label{ex:bands}
Fix \(\eps_0>0\), sufficiently small for the later packing
argument, and any fixed loss exponent \(C_*\). For large \(T\),
there are \(Q\) ordered disjoint bands
\([B_i,2B_i]\subset(2a_0,T/2)\), independent of the contact,
such that
\begin{equation}
 B_1^{1/10}\gg F_Q(T)^3,\qquad
 S_h(a_0)^{C_*}\ll B_1^{\eps_0},\qquad
 S_1(2B_i)^{C_*}\ll B_{i+1}^{\eps_0}.
 \label{ex:band-spacing}
\end{equation}
For \(Q=1\), only the first condition is needed, and a single
band of size comparable to \(T\) suffices. Every contact path has
a traversal contained in one band, from \(a=B\) to \(a=2B\),
which contains no merger and on which \(k\ge B^{-\eps_0}\).
The estimates \eqref{ex:bank}, \eqref{ex:dilations},
\eqref{ex:derivatives}, \eqref{ex:weighted-scalar}, and
\eqref{ex:comparison}--\eqref{ex:transition} hold on that traversal.
If it has several wells, the bound
\(\dot Z\le a^{-P_0}I\) holds on every preceding stage with
several wells as well, and the initial value of \(Z\) is the
arbitrarily small quantity in \Cref{ex:initialization}.
\end{lemma}

\begin{proof}
We spell out the rate comparison to keep track of the inverse.
Let \(Q\ge2\) and \(j=2Q-2\). Above the iterated-log threshold,
\begin{equation}
 F_{Q-1}^{-1}(u)
   =\exp_j\bigl((\log_j u)^2\bigr).
 \label{ex:exact-inverse}
\end{equation}
For \(b\) chosen as above, inversion of any fixed power of
\(S_1(x)=e^{Cb(x)}\) first takes a logarithm, then applies
\(F_{Q-1}\), and then inverts \((\log x)^3\). In particular
the inverse retains the outer exponential from this last step.
For example when \(Q=2\) its leading expression, up to fixed
losses, is
\[
 \exp\!\left(\bigl[F_1(c\log u)\bigr]^{1/3}\right).
\]
It follows directly from \eqref{ex:exact-inverse} that, for some
fixed \(\varsigma>0\), every such inverse dominates
\[
 \exp_{2Q-1}\bigl((\log_{2Q-1}u)^{\varsigma}\bigr).
\]
Any fixed number of inverse compositions, including fixed powers
of the argument and value, retains this lower bound with a smaller
positive \(\varsigma\). To compare with \(F_Q(T)^C\), take
\(2Q-1\) logarithms: the lower bound has last value
\(\exp(\varsigma\log_{2Q}T)\), whereas the latter has last value
\(\exp(\sqrt{\log_{2Q}T})\), up to lower order terms.
The former is larger for every fixed \(C\).

The same comparison shows \(b(T)=o(F_Q(T)^c)\) for each fixed
\(c>0\), and hence \(b(T)+b(a_0)=o(a_0)\), as used above.
It also shows that after a fixed number of the inverse operations,
\(\log B_1\gg b(a_0)\). This stronger comparison absorbs the
initial factor \(S_h(a_0)\), including any fixed power of it.
Enlarge \(C_*\), if necessary, to exceed the finitely many powers
arising from \(Q-1\) applications of \eqref{ex:merge-k}; proving the
stronger separation also proves the originally requested one.
Now choose \(B_Q\) to be a small fixed multiple of \(T\), and
choose the earlier bands successively by inverse operations with
enough fixed slack to ensure \eqref{ex:band-spacing}.
The comparisons just proved give all the lower-end requirements.
For \(Q=1\), \(F_1(T)=T^{o(1)}\), so the claimed single band
works and no inverse construction is required.

Consider the path from initialization to its first hit of each
successive level \(2B_i\). There are \(Q\) resulting consecutive
time intervals and fewer than \(Q\) merger events, counting a
common-well initialization if needed. At least one interval has no
event. On it take the last preceding hit of \(B_i\) before the
first hit of \(2B_i\); this is a traversal staying in that band.
All earlier mergers have levels at most \(2B_{i-1}\) (or are
initialization events). Iterating \eqref{ex:merge-k}, with at most
\(Q-1\) steps, bounds their total loss by a fixed power of
\(S_1(2B_{i-1})\), the initial \(S_h(a_0)\), and fixed constants.
Equation \eqref{ex:band-spacing} and the explicit solution of
\eqref{ex:k-evolution} then give \(k\ge B_i^{-\eps_0}\)
throughout the traversal. The remaining assertions are the estimates
already proved. The number of wells only decreases, so a traversal
with several wells has had several wells at every earlier stage.
\end{proof}
